\ifdefined\inMainDoc\else
\documentclass[12pt,a4paper]{article}
% ================================================================
%  PRÉAMBULE COMMUN - ANNALES CORRIGÉES
%  Université de Kara - FaST - Licence Fondamentale MPC
%  Design optimisé impression noir et blanc
% ================================================================

% -- ENCODAGE & LANGUE --
\usepackage[utf8]{inputenc}
\usepackage[T1]{fontenc}
\usepackage[french]{babel}
\usepackage{lmodern}
\usepackage{microtype}

% -- MISE EN PAGE --
\usepackage[a4paper, top=2.8cm, bottom=2.5cm, left=2.5cm, right=2.5cm, headheight=14pt]{geometry}
\usepackage{setspace}
\setstretch{1.2}
\usepackage{parskip}
\setlength{\parindent}{1.2em}
\setlength{\parskip}{4pt}

% -- MATHS --
\usepackage{amsmath, amssymb, mathtools, amsthm}
\usepackage{mathrsfs}
\usepackage{dsfont}
\usepackage{bm}
\usepackage{cancel}
\usepackage{textcomp}
% Définition de \oiint (intégrale de surface fermée) sans package supplémentaire
\newcommand{\oiint}{\mathop{\iint\!\!\!\!\!\!\!\bigcirc}\limits}

% -- PHYSIQUE & CHIMIE --
\usepackage{physics}
\usepackage{siunitx}
% Lorsque physics est chargé, siunitx omet \qty. On le restaure ici :
\AtBeginDocument{\RenewCommandCopy\qty\SI}
\sisetup{locale = FR, detect-all, output-decimal-marker={,}}
% Commande de substitution pour les unités posant problème avec pdflatex
% (ohm, degreeCelsius, micro, angstrom, percent utilisent des caractères non-ASCII)
\newcommand{\qtyohm}[1]{\ensuremath{\num{#1}\,\Omega}}
\newcommand{\qtydegc}[1]{\ensuremath{\num{#1}\,{}^\circ\text{C}}}
\newcommand{\qtyangstrom}[1]{\ensuremath{\num{#1}\,\text{\AA}}}
\newcommand{\qtypercent}[1]{\ensuremath{\num{#1}\,\%}}
\newcommand{\qtyum}[1]{\ensuremath{\num{#1}\,\mu\text{m}}}
\usepackage[version=4]{mhchem}

% -- GRAPHISMES --
\usepackage{graphicx}
\usepackage{float}
\usepackage{caption}
\usepackage{subcaption}
\usepackage{wrapfig}
\usepackage{tikz}
\usetikzlibrary{calc, arrows.meta, patterns, shapes.geometric}
\usepackage{pgfplots}
\pgfplotsset{compat=1.18}

% -- TABLEAUX --
\usepackage{booktabs}
\usepackage{array}
\usepackage{tabularx}
\usepackage{multirow}

% -- LISTES --
\usepackage{enumitem}
\setlist[enumerate]{topsep=3pt, itemsep=2pt, parsep=0pt}
\setlist[itemize]{topsep=3pt, itemsep=2pt, parsep=0pt, label={.}}

% -- BOÎTES (noir et blanc) --
\usepackage{mdframed}
\usepackage{tcolorbox}
\tcbuselibrary{skins, breakable, theorems}

% Boîte EXERCICE : encadré avec filet épais, fond blanc
\newmdenv[
  linewidth=1.2pt,
  linecolor=black,
  roundcorner=3pt,
  skipabove=10pt,
  skipbelow=6pt,
  innertopmargin=8pt,
  innerbottommargin=8pt,
  innerleftmargin=10pt,
  innerrightmargin=10pt,
  backgroundcolor=white,
  frametitle={\bfseries\small Exercice},
  frametitlealignment=\raggedright,
  frametitleaboveskip=4pt,
  frametitlebelowskip=4pt,
]{boiteexercice}

% Boîte CORRIGÉ : fond gris très clair + filet fin
\newmdenv[
  linewidth=0.6pt,
  linecolor=black,
  leftline=true,
  rightline=false,
  topline=false,
  bottomline=false,
  linecolor=black,
  innerleftmargin=12pt,
  innerrightmargin=8pt,
  innertopmargin=6pt,
  innerbottommargin=6pt,
  backgroundcolor=gray!8,
  skipabove=4pt,
  skipbelow=8pt,
]{boitecorrige}

% Boîte RÉSUMÉ DE COURS : double encadré
\newmdenv[
  linewidth=0.8pt,
  linecolor=black,
  roundcorner=2pt,
  backgroundcolor=gray!12,
  innertopmargin=10pt,
  innerbottommargin=10pt,
  innerleftmargin=12pt,
  innerrightmargin=12pt,
  skipabove=12pt,
  skipbelow=8pt,
]{boiteresume}

% Boîte ATTENTION/REMARQUE : barre gauche épaisse
\newmdenv[
  leftline=true,
  rightline=false,
  topline=false,
  bottomline=false,
  linewidth=3pt,
  linecolor=black,
  backgroundcolor=gray!5,
  innerleftmargin=12pt,
  innerrightmargin=8pt,
  innertopmargin=5pt,
  innerbottommargin=5pt,
  skipabove=6pt,
  skipbelow=6pt,
]{boiteremarque}

% -- DIFFICULTÉ DES EXERCICES --
\newcommand{\facile}{$\star$}
\newcommand{\moyen}{$\star\!\star$}
\newcommand{\difficile}{$\star\!\star\!\star$}
\newcommand{\niveau}[1]{\hfill\textbf{#1}\hspace{2pt}}

% -- EN-TÊTES ET PIEDS DE PAGE --
\usepackage{fancyhdr}
\pagestyle{fancy}
\fancyhf{}
\renewcommand{\headrulewidth}{0.5pt}
\renewcommand{\footrulewidth}{0.3pt}

% -- HYPERLIENS --
\usepackage[hidelinks]{hyperref}
\hypersetup{
    colorlinks=false,
    pdftitle={Annale Corrigée - Licence MPC - Université de Kara}
}

% -- TITRES --
\usepackage{titlesec}
\titleformat{\section}{\Large\bfseries}{\thesection.}{0.8em}{}[\vspace{2pt}\hrule\vspace{4pt}]
\titleformat{\subsection}{\large\bfseries}{\thesubsection.}{0.7em}{}
\titleformat{\subsubsection}{\normalsize\bfseries}{\thesubsubsection.}{0.6em}{}

% -- THÉORÈMES --
\theoremstyle{definition}
\newtheorem{defn}{Définition}[section]
\newtheorem{prop}{Propriété}[section]
\newtheorem{thm}{Théorème}[section]
\newtheorem{corol}{Corollaire}[section]
\newtheorem{exemple}{Exemple}[section]

% -- COMMANDES MATHÉMATIQUES UTILES --
\newcommand{\vect}[1]{\overrightarrow{#1}}
\newcommand{\R}{\mathbb{R}}
\newcommand{\N}{\mathbb{N}}
\newcommand{\Z}{\mathbb{Z}}
\newcommand{\C}{\mathbb{C}}
\renewcommand{\d}{\mathrm{d}}
\newcommand{\deriv}[2]{\dfrac{\d #1}{\d #2}}
\newcommand{\dpart}[2]{\dfrac{\partial #1}{\partial #2}}
\newcommand{\rot}{\overrightarrow{\mathrm{rot}}\,}
\renewcommand{\div}{\mathrm{div}\,}
\renewcommand{\grad}{\overrightarrow{\mathrm{grad}}\,}
\newcommand{\e}{\mathrm{e}}
\newcommand{\ii}{\mathrm{i}}
\renewcommand{\Tr}{\mathrm{Tr}}

% Numérotation des équations par section
\numberwithin{equation}{section}

% -- RÉSULTATS ENCADRÉS --
\newcommand{\res}[1]{\[\boxed{#1}\]}

% -- REMARQUE INLINE --
\newcommand{\rmq}[1]{\begin{boiteremarque}\textbf{Remarque.} #1\end{boiteremarque}}

% -- MISE EN VALEUR DU RÉSULTAT FINAL --
\newcommand{\reponse}[1]{%
  \begin{center}
    \fbox{\begin{minipage}{0.92\textwidth}\centering #1\end{minipage}}
  \end{center}%
}


% En-têtes et pieds de page
\lhead{\textbf{Analyse dans $\mathbb{R}^n$} - Licence 2}
\rhead{Université de Kara - FaST}
\cfoot{\thepage}
\lfoot{\textit{Licence Fondamentale MPC}}
\rfoot{\textit{Annales corrigées}}

\begin{document}
\fi

\begin{center}
  {\LARGE\bfseries Analyse dans $\mathbb{R}^n$}\\[6pt]
  {\large Licence 2 - Mathématiques, Physique, Chimie}\\[4pt]
  {\normalsize Université de Kara - Faculté des Sciences et Techniques}
\end{center}

\vspace{8pt}
\hrule
\vspace{12pt}

% ================================================================
\section{Résumé de cours}
% ================================================================

\begin{boiteresume}
\textbf{1. Fonctions de plusieurs variables}

Une fonction $f : U \subset \mathbb{R}^n \to \mathbb{R}$ est \textbf{continue en $a$} si
$\displaystyle\lim_{x \to a} f(x) = f(a)$,
ce qui, pour $n \geq 2$, doit être vérifié quelle que soit la direction d'approche.

\textbf{Dérivées partielles} :
\[
\dpart{f}{x_i}(a) = \lim_{h \to 0} \frac{f(a + h\,e_i) - f(a)}{h},
\]
où $e_i$ est le $i$-ème vecteur de la base canonique.

\medskip
\textbf{2. Différentielle totale et gradient}

$f$ est \textbf{différentiable en $a$} s'il existe une forme linéaire $L$ telle que :
\[
f(a+h) = f(a) + L(h) + o(\|h\|).
\]
La différentielle totale est alors $\mathrm{d}f_a = \displaystyle\sum_{i=1}^n \dpart{f}{x_i}(a)\,\mathrm{d}x_i$.

Le \textbf{gradient} de $f$ en $a$ est :
\[
\overrightarrow{\nabla f}(a) = \sum_{i=1}^n \dpart{f}{x_i}(a)\,\vec{e}_i.
\]

\textbf{Règle de dérivation en chaîne} : si $x(t) : \mathbb{R} \to \mathbb{R}^n$ est dérivable et $f$ est différentiable,
\[
\frac{\mathrm{d}}{\mathrm{d}t}\bigl[f(x(t))\bigr] = \overrightarrow{\nabla f}(x(t)) \cdot x'(t) = \sum_{i=1}^n \dpart{f}{x_i}\,\dot{x}_i.
\]

Pour $g(s,t) = f(x(s,t), y(s,t))$ :
\[
\dpart{g}{s} = \dpart{f}{x}\,\dpart{x}{s} + \dpart{f}{y}\,\dpart{y}{s}.
\]

\medskip
\textbf{3. Extrema libres et liés}

\textbf{Extremum libre} : si $f$ admet un extremum local en $a$, alors $\overrightarrow{\nabla f}(a) = \vec{0}$ (point critique). La nature est déterminée par la matrice hessienne $H_f$ :
\[
H_f(a) = \left(\frac{\partial^2 f}{\partial x_i \partial x_j}(a)\right)_{i,j}.
\]
En deux variables :
\begin{itemize}[nosep]
  \item si $\det H_f > 0$ et $\partial^2 f/\partial x^2 > 0$, c'est un \textbf{minimum} ;
  \item si $\det H_f > 0$ et $\partial^2 f/\partial x^2 < 0$, c'est un \textbf{maximum} ;
  \item si $\det H_f < 0$, c'est un \textbf{point selle}.
\end{itemize}

\textbf{Multiplicateurs de Lagrange.} Pour optimiser $f(x,y)$ sous $g(x,y) = 0$, on cherche $(x,y,\lambda)$ vérifiant :
\[
\overrightarrow{\nabla f}(x,y) = \lambda\,\overrightarrow{\nabla g}(x,y) \quad \text{et} \quad g(x,y) = 0.
\]

\medskip
\textbf{4. Intégrales doubles}

\[
\iint_D f(x,y)\,\mathrm{d}x\,\mathrm{d}y = \int_{a}^{b}\!\left(\int_{\varphi_1(x)}^{\varphi_2(x)} f(x,y)\,\mathrm{d}y\right)\!\mathrm{d}x \quad \text{(théorème de Fubini).}
\]

\textbf{Changement de variables} : $(x,y) \to (u,v)$ de jacobien $J = \dpart{(x,y)}{(u,v)}$,
\[
\iint_D f(x,y)\,\mathrm{d}x\,\mathrm{d}y = \iint_{D'} f\bigl(x(u,v), y(u,v)\bigr)\,|J|\,\mathrm{d}u\,\mathrm{d}v.
\]

\textbf{Coordonnées polaires} : $x = r\cos\theta$, $y = r\sin\theta$, $|J| = r$.

\medskip
\textbf{5. Intégrales triples et coordonnées curvilignes}

\[
\iiint_V f\,\mathrm{d}V = \iiint f\,|J|\,\mathrm{d}u\,\mathrm{d}v\,\mathrm{d}w.
\]

\textbf{Coordonnées cylindriques} : $x = r\cos\theta$, $y = r\sin\theta$, $z = z$, $|J| = r$.

\textbf{Coordonnées sphériques} : $x = r\sin\phi\cos\theta$, $y = r\sin\phi\sin\theta$, $z = r\cos\phi$, $|J| = r^2\sin\phi$.
\end{boiteresume}

% ================================================================
\section{Exercices corrigés}
% ================================================================

% ----------------------
\begin{boiteexercice}
\textbf{Exercice 1 - Continuité et dérivées partielles}\niveau{\facile}

Soit $f : \mathbb{R}^2 \to \mathbb{R}$ définie par :
\[
f(x,y) = \begin{cases} \dfrac{x^2 y}{x^2 + y^2} & \text{si } (x,y) \neq (0,0), \\ 0 & \text{si } (x,y) = (0,0). \end{cases}
\]

\begin{enumerate}
  \item Calculer $\dpart{f}{x}(0,0)$ et $\dpart{f}{y}(0,0)$.
  \item $f$ est-elle continue en $(0,0)$ ?
  \item $f$ est-elle différentiable en $(0,0)$ ?
\end{enumerate}
\end{boiteexercice}

\begin{boitecorrige}
\textbf{1. Dérivées partielles en $(0,0)$.}

Par définition :
\[
\dpart{f}{x}(0,0) = \lim_{h \to 0} \frac{f(h,0) - f(0,0)}{h} = \lim_{h \to 0} \frac{0}{h} = 0.
\]
\[
\dpart{f}{y}(0,0) = \lim_{k \to 0} \frac{f(0,k) - f(0,0)}{k} = \lim_{k \to 0} \frac{0}{k} = 0.
\]

\textbf{2. Continuité en $(0,0)$.}

On approche $(0,0)$ le long de la droite $y = tx$ ($t \in \mathbb{R}$, $t \neq 0$) :
\[
f(x, tx) = \frac{x^2 \cdot tx}{x^2 + t^2x^2} = \frac{tx}{1 + t^2} \xrightarrow[x \to 0]{} 0.
\]
Avec $t = 1$ (droite $y=x$) : $f(x,x) = \dfrac{x}{2} \to 0$.

Vérifions avec $y = x^2$ : $f(x, x^2) = \dfrac{x^2 \cdot x^2}{x^2 + x^4} = \dfrac{x^2}{1+x^2} \to 0$.

Toutes les limites tendent vers 0 et $f(0,0) = 0$. La fonction $f$ est continue en $(0,0)$.

\textbf{3. Différentiabilité en $(0,0)$.}

Si $f$ était différentiable en $(0,0)$, avec $\mathrm{d}f_0 = 0\,\mathrm{d}x + 0\,\mathrm{d}y = 0$, il faudrait :
\[
\lim_{(h,k)\to(0,0)} \frac{f(h,k) - 0}{\sqrt{h^2+k^2}} = 0.
\]
Avec $h = r\cos\theta$, $k = r\sin\theta$ :
\[
\frac{f(h,k)}{\sqrt{h^2+k^2}} = \frac{r^2\cos^2\theta \cdot r\sin\theta}{r^2 \cdot r} = \cos^2\theta\sin\theta.
\]
Cette expression dépend de $\theta$ et ne tend pas vers $0$ uniformément. Par exemple, pour $\theta = \pi/4$, elle vaut $1/2\sqrt{2} \neq 0$. Donc $f$ n'est pas différentiable en $(0,0)$.
\end{boitecorrige}

% ----------------------
\begin{boiteexercice}
\textbf{Exercice 2 - Règle de dérivation en chaîne}\niveau{\facile}

Soit $f(x,y) = x^2 e^y + y\sin(x)$. On pose $x(t) = \cos t$ et $y(t) = t^2$.

\begin{enumerate}
  \item Calculer $\overrightarrow{\nabla f}(x,y)$.
  \item En utilisant la règle de dérivation en chaîne, calculer $\dfrac{\mathrm{d}}{\mathrm{d}t}\bigl[f(x(t), y(t))\bigr]$.
  \item Vérifier le résultat en développant directement $g(t) = f(\cos t, t^2)$ puis en dérivant.
\end{enumerate}
\end{boiteexercice}

\begin{boitecorrige}
\textbf{1. Gradient de $f$.}

\[
\dpart{f}{x} = 2x\,\mathrm{e}^y + y\cos x, \qquad \dpart{f}{y} = x^2\,\mathrm{e}^y + \sin x.
\]
\[
\overrightarrow{\nabla f}(x,y) = \bigl(2x\,\mathrm{e}^y + y\cos x\bigr)\,\vec{i} + \bigl(x^2\,\mathrm{e}^y + \sin x\bigr)\,\vec{j}.
\]

\textbf{2. Dérivée en chaîne.}

$\dot{x}(t) = -\sin t$, $\dot{y}(t) = 2t$.

\[
\frac{\mathrm{d}g}{\mathrm{d}t} = \dpart{f}{x}\,\dot{x} + \dpart{f}{y}\,\dot{y}
= \bigl(2\cos t\,\mathrm{e}^{t^2} + t^2\cos(\cos t)\bigr)(-\sin t) + \bigl(\cos^2 t\,\mathrm{e}^{t^2} + \sin(\cos t)\bigr)\cdot 2t.
\]

En développant :
\[
\frac{\mathrm{d}g}{\mathrm{d}t} = -2\sin t\cos t\,\mathrm{e}^{t^2} - t^2\sin t\cos(\cos t) + 2t\cos^2 t\,\mathrm{e}^{t^2} + 2t\sin(\cos t).
\]

\textbf{3. Vérification directe.}

$g(t) = \cos^2 t\,\mathrm{e}^{t^2} + t^2\sin(\cos t)$.

$\dfrac{\mathrm{d}}{\mathrm{d}t}\bigl[\cos^2 t\,\mathrm{e}^{t^2}\bigr] = -2\sin t\cos t\,\mathrm{e}^{t^2} + 2t\cos^2 t\,\mathrm{e}^{t^2}$.

$\dfrac{\mathrm{d}}{\mathrm{d}t}\bigl[t^2\sin(\cos t)\bigr] = 2t\sin(\cos t) + t^2\cos(\cos t)\cdot(-\sin t)$.

On retrouve exactement le même résultat.
\end{boitecorrige}

% ----------------------
\begin{boiteexercice}
\textbf{Exercice 3 - Extrema libres et liés}\niveau{\moyen}

\textbf{Partie A.} Chercher les extrema locaux de $f(x,y) = x^3 + y^3 - 3xy$.

\textbf{Partie B.} Maximiser $f(x,y) = xy$ sur l'ellipse $\dfrac{x^2}{4} + y^2 = 1$ en utilisant les multiplicateurs de Lagrange.
\end{boiteexercice}

\begin{boitecorrige}
\textbf{Partie A : extrema libres.}

Points critiques : $\overrightarrow{\nabla f} = \vec{0}$.
\[
\dpart{f}{x} = 3x^2 - 3y = 0 \Rightarrow y = x^2.
\qquad
\dpart{f}{y} = 3y^2 - 3x = 0 \Rightarrow x = y^2.
\]
En substituant $y = x^2$ dans $x = y^2 = x^4$, on obtient $x^4 - x = 0$, soit $x(x^3-1) = 0$. Donc $x = 0$ ou $x = 1$.

Points critiques : $(0,0)$ et $(1,1)$.

Matrice hessienne :
\[
H_f(x,y) = \begin{pmatrix} 6x & -3 \\ -3 & 6y \end{pmatrix}.
\]

En $(0,0)$ : $\det H_f = 0 \times 0 - 9 = -9 < 0$ : c'est un \textbf{point selle} (ni maximum ni minimum).

En $(1,1)$ : $\det H_f = 6\times6 - 9 = 27 > 0$ et $H_{11} = 6 > 0$ : c'est un \textbf{minimum local}. $f(1,1) = 1 + 1 - 3 = -1$.

\textbf{Partie B : multiplicateurs de Lagrange.}

Contrainte : $g(x,y) = \dfrac{x^2}{4} + y^2 - 1 = 0$.

\[
\overrightarrow{\nabla f} = \lambda\,\overrightarrow{\nabla g} \Leftrightarrow
\begin{cases} y = \lambda\,\dfrac{x}{2} \\[4pt] x = 2\lambda\,y \end{cases}.
\]

En substituant la première dans la seconde : $x = 2\lambda \cdot \dfrac{\lambda x}{2} = \lambda^2 x$.

Si $x \neq 0$ : $\lambda^2 = 1$, soit $\lambda = \pm 1$.

Pour $\lambda = 1$ : $y = x/2$. Contrainte : $\dfrac{x^2}{4} + \dfrac{x^2}{4} = 1 \Rightarrow x^2 = 2 \Rightarrow x = \pm\sqrt{2}$, $y = \pm\dfrac{\sqrt{2}}{2}$.

Pour $\lambda = -1$ : $y = -x/2$. Mêmes valeurs de $|x|$, mais $y = -x/2$.

Si $x = 0$ : $y^2 = 1$, $f(0,\pm1) = 0$.

Les valeurs candidates :
\[
f\!\left(\sqrt{2},\frac{\sqrt{2}}{2}\right) = \sqrt{2}\cdot\frac{\sqrt{2}}{2} = 1, \quad
f\!\left(-\sqrt{2},-\frac{\sqrt{2}}{2}\right) = 1, \quad
f\!\left(\sqrt{2},-\frac{\sqrt{2}}{2}\right) = -1.
\]

Le \textbf{maximum} de $xy$ sur l'ellipse est $\boxed{1}$, atteint en $\bigl(\sqrt{2}, \tfrac{\sqrt{2}}{2}\bigr)$ et $\bigl(-\sqrt{2}, -\tfrac{\sqrt{2}}{2}\bigr)$.
\end{boitecorrige}

% ----------------------
\begin{boiteexercice}
\textbf{Exercice 4 - Intégrale double, changement en polaires}\niveau{\moyen}

\textbf{Partie A.} Calculer $\displaystyle I = \iint_D (x^2 + y^2)\,\mathrm{d}x\,\mathrm{d}y$ où $D$ est le disque de rayon $R$ centré à l'origine.

\textbf{Partie B.} Calculer $\displaystyle J = \int_0^{+\infty} \mathrm{e}^{-x^2}\,\mathrm{d}x$ en utilisant $I_0 = \iint_{\mathbb{R}^2} \mathrm{e}^{-(x^2+y^2)}\,\mathrm{d}x\,\mathrm{d}y$.
\end{boiteexercice}

\begin{boitecorrige}
\textbf{Partie A.} On passe en coordonnées polaires : $x = r\cos\theta$, $y = r\sin\theta$, $\mathrm{d}x\,\mathrm{d}y = r\,\mathrm{d}r\,\mathrm{d}\theta$.

$D$ correspond à $r \in [0, R]$, $\theta \in [0, 2\pi]$.

\[
I = \int_0^{2\pi}\!\int_0^R r^2 \cdot r\,\mathrm{d}r\,\mathrm{d}\theta
= \int_0^{2\pi}\mathrm{d}\theta \cdot \int_0^R r^3\,\mathrm{d}r
= 2\pi \cdot \frac{R^4}{4} = \frac{\pi R^4}{2}.
\]

\textbf{Partie B.}

\[
I_0 = \int_{-\infty}^{+\infty}\!\int_{-\infty}^{+\infty} \mathrm{e}^{-(x^2+y^2)}\,\mathrm{d}x\,\mathrm{d}y
= \left(\int_{-\infty}^{+\infty} \mathrm{e}^{-x^2}\,\mathrm{d}x\right)^2 = (2J)^2 = 4J^2.
\]

En passant en polaires sur $\mathbb{R}^2$ :
\[
I_0 = \int_0^{2\pi}\!\int_0^{+\infty} \mathrm{e}^{-r^2}\,r\,\mathrm{d}r\,\mathrm{d}\theta
= 2\pi \int_0^{+\infty} r\,\mathrm{e}^{-r^2}\,\mathrm{d}r.
\]
Avec le changement $u = r^2$ : $\displaystyle\int_0^{+\infty} r\,\mathrm{e}^{-r^2}\,\mathrm{d}r = \dfrac{1}{2}$.

Donc $I_0 = \pi$, ce qui donne $4J^2 = \pi$, soit :
\[
\boxed{J = \int_0^{+\infty} \mathrm{e}^{-x^2}\,\mathrm{d}x = \frac{\sqrt{\pi}}{2}.}
\]
\end{boitecorrige}

% ----------------------
\begin{boiteexercice}
\textbf{Exercice 5 - Intégrale triple en coordonnées cylindriques et sphériques}\niveau{\moyen}

\textbf{Partie A.} Calculer le volume du cylindre d'axe $Oz$, de rayon $a$, compris entre les plans $z = 0$ et $z = h$ :
\[
V_1 = \iiint_C \mathrm{d}V.
\]

\textbf{Partie B.} Calculer l'intégrale
\[
V_2 = \iiint_B \bigl(x^2 + y^2 + z^2\bigr)\,\mathrm{d}V,
\]
où $B$ est la boule de rayon $R$ centrée à l'origine.
\end{boiteexercice}

\begin{boitecorrige}
\textbf{Partie A.} En coordonnées cylindriques $(r, \theta, z)$ avec $r \in [0,a]$, $\theta \in [0, 2\pi]$, $z \in [0, h]$ et $\mathrm{d}V = r\,\mathrm{d}r\,\mathrm{d}\theta\,\mathrm{d}z$ :

\[
V_1 = \int_0^{2\pi}\!\int_0^a\!\int_0^h r\,\mathrm{d}z\,\mathrm{d}r\,\mathrm{d}\theta
= 2\pi \cdot \frac{a^2}{2} \cdot h = \pi a^2 h.
\]
On retrouve bien la formule classique du volume d'un cylindre.

\textbf{Partie B.} En coordonnées sphériques $(r, \phi, \theta)$ avec $r \in [0,R]$, $\phi \in [0, \pi]$, $\theta \in [0, 2\pi]$ et $\mathrm{d}V = r^2\sin\phi\,\mathrm{d}r\,\mathrm{d}\phi\,\mathrm{d}\theta$ :

$x^2 + y^2 + z^2 = r^2$.

\[
V_2 = \int_0^{2\pi}\!\int_0^{\pi}\!\int_0^R r^2 \cdot r^2\sin\phi\,\mathrm{d}r\,\mathrm{d}\phi\,\mathrm{d}\theta
= \left(\int_0^{2\pi}\mathrm{d}\theta\right)\!\left(\int_0^{\pi}\sin\phi\,\mathrm{d}\phi\right)\!\left(\int_0^R r^4\,\mathrm{d}r\right).
\]

\[
= 2\pi \cdot \bigl[-\cos\phi\bigr]_0^{\pi} \cdot \frac{R^5}{5}
= 2\pi \cdot 2 \cdot \frac{R^5}{5} = \frac{4\pi R^5}{5}.
\]
\end{boitecorrige}

% ----------------------
\begin{boiteexercice}
\textbf{Exercice 6 - Changement de variables, jacobien}\niveau{\difficile}

On considère l'intégrale
\[
I = \iint_D \frac{y}{x}\,\mathrm{d}x\,\mathrm{d}y,
\]
où $D$ est le domaine délimité par les courbes $y = x$, $y = 2x$, $xy = 1$ et $xy = 2$ (avec $x > 0$).

On pose $u = \dfrac{y}{x}$ et $v = xy$.

\begin{enumerate}
  \item Déterminer le domaine $D'$ dans le plan $(u,v)$.
  \item Calculer le jacobien $\dpart{(x,y)}{(u,v)}$.
  \item Calculer $I$.
\end{enumerate}
\end{boiteexercice}

\begin{boitecorrige}
\textbf{1. Domaine $D'$ dans le plan $(u,v)$.}

Les courbes frontières :
\begin{itemize}
  \item $y = x$ donne $u = 1$ ; $y = 2x$ donne $u = 2$.
  \item $xy = 1$ donne $v = 1$ ; $xy = 2$ donne $v = 2$.
\end{itemize}
Donc $D'$ est le rectangle $[1,2] \times [1,2]$ dans le plan $(u,v)$.

\textbf{2. Jacobien.}

On exprime $x$ et $y$ en fonction de $u$ et $v$ : $y = ux$ et $v = xy = ux^2$, donc $x = \sqrt{v/u}$ et $y = \sqrt{uv}$.

\[
\dpart{x}{u} = -\frac{1}{2}\sqrt{\frac{v}{u^3}}, \quad
\dpart{x}{v} = \frac{1}{2}\sqrt{\frac{1}{uv}}, \quad
\dpart{y}{u} = \frac{1}{2}\sqrt{\frac{v}{u}}, \quad
\dpart{y}{v} = \frac{1}{2}\sqrt{\frac{u}{v}}.
\]

\[
J = \dpart{(x,y)}{(u,v)} = \dpart{x}{u}\dpart{y}{v} - \dpart{x}{v}\dpart{y}{u}
= -\frac{1}{4u} - \frac{1}{4u} = -\frac{1}{2u}.
\]

Donc $|J| = \dfrac{1}{2u}$.

\textbf{3. Calcul de $I$.}

\[
I = \iint_{D'} u \cdot \frac{1}{2u}\,\mathrm{d}u\,\mathrm{d}v
= \iint_{D'} \frac{1}{2}\,\mathrm{d}u\,\mathrm{d}v
= \frac{1}{2} \int_1^2\!\int_1^2 \mathrm{d}u\,\mathrm{d}v
= \frac{1}{2} \times 1 \times 1 = \frac{1}{2}.
\]
\end{boitecorrige}

% ----------------------
\begin{boiteexercice}
\textbf{Exercice 7 - Continuité en plusieurs variables : définition $\varepsilon$-$\delta$}\niveau{\moyen}

Soit $f : \mathbb{R}^2 \to \mathbb{R}$ définie par :
\[
f(x,y) = \begin{cases}\dfrac{xy}{x^2+y^2} & \text{si }(x,y)\neq(0,0)\\0 & \text{si }(x,y)=(0,0)\end{cases}
\]
\begin{enumerate}
  \item Calculer la limite de $f(x,y)$ quand $(x,y)\to(0,0)$ le long de l'axe $Ox$ ($y=0$) et le long de la droite $y=x$.
  \item En déduire que $f$ est \textbf{discontinue} en $(0,0)$.
  \item Montrer que $|f(x,y)|\leq 1/2$ pour tout $(x,y)\neq(0,0)$ (inégalité arithmético-géométrique).
  \item Calculer les dérivées partielles $\partial f/\partial x(0,0)$ et $\partial f/\partial y(0,0)$.
  \item $f$ est-elle différentiable en $(0,0)$ ? Justifier.
\end{enumerate}
\end{boiteexercice}

\begin{boitecorrige}
\textbf{1.} Le long de $y=0$ ($x\neq0$) : $f(x,0)=0 \to 0$. Le long de $y=x$ ($x\neq0$) : $f(x,x)=\dfrac{x^2}{2x^2}=\dfrac{1}{2}\to\dfrac{1}{2}$.

\textbf{2.} Les deux limites sont différentes ($0\neq 1/2$) : la limite de $f$ en $(0,0)$ n'existe pas. Donc $f$ est \textbf{discontinue} en $(0,0)$.

\textbf{3.} Par l'inégalité arithmético-géométrique : $x^2 + y^2 \geq 2|xy|$, donc :
\[
|f(x,y)| = \frac{|xy|}{x^2+y^2} \leq \frac{|xy|}{2|xy|} = \frac{1}{2}
\]

\textbf{4.} $\dpart{f}{x}(0,0) = \lim_{h\to0}\dfrac{f(h,0)-f(0,0)}{h} = \lim_{h\to0}\dfrac{0}{h} = 0$. De même $\dpart{f}{y}(0,0) = 0$.

\textbf{5.} Si $f$ était différentiable en $(0,0)$, avec $\mathrm{d}f_{(0,0)} = 0$, il faudrait $f(h,k)/\sqrt{h^2+k^2}\to0$. En prenant $h = r\cos\theta$, $k = r\sin\theta$ :
\[
\frac{f(h,k)}{\sqrt{h^2+k^2}} = \frac{r^2\cos\theta\sin\theta}{r^2} = \cos\theta\sin\theta\frac{1}{2}\sin(2\theta)
\]
Cette expression dépend de $\theta$ et ne tend pas vers $0$ uniformément. $f$ n'est \textbf{pas différentiable} en $(0,0)$. Cet exemple montre que l'existence des dérivées partielles ne suffit pas à garantir la différentiabilité.
\end{boitecorrige}

% ----------------------
\begin{boiteexercice}
\textbf{Exercice 8 - Matrice jacobienne et applications}\niveau{\moyen}

\begin{enumerate}
  \item Soit $F : \mathbb{R}^2\to\mathbb{R}^2$ définie par $F(x,y) = (x^2+y, xy+y^2)$. Calculer la matrice jacobienne $J_F(x,y)$.
  \item Calculer $J_F(1,2)$ et $\det J_F(1,2)$.
  \item Soit la transformation polaire $\Phi : (r,\theta)\mapsto(r\cos\theta, r\sin\theta)$. Calculer $J_\Phi$ et $\det J_\Phi$.
  \item En utilisant la règle de la chaîne, calculer $J_{G\circ F}$ en $(1,2)$ où $G(u,v) = (u^2-v, uv)$.
\end{enumerate}
\end{boiteexercice}

\begin{boitecorrige}
\textbf{1.} $F = (F_1, F_2)$ avec $F_1 = x^2+y$ et $F_2 = xy+y^2$ :
\[
J_F(x,y) = \begin{pmatrix}\partial F_1/\partial x & \partial F_1/\partial y\\\partial F_2/\partial x & \partial F_2/\partial y\end{pmatrix} = \begin{pmatrix}2x & 1\\y & x+2y\end{pmatrix}
\]

\textbf{2.} En $(1,2)$ :
\[
J_F(1,2) = \begin{pmatrix}2 & 1\\2 & 1+4\end{pmatrix} = \begin{pmatrix}2 & 1\\2 & 5\end{pmatrix} \qquad \det J_F(1,2) = 2\times5 - 1\times2 = 8
\]
Comme $\det J_F \neq 0$, $F$ est un difféomorphisme local en $(1,2)$.

\textbf{3.} $x = r\cos\theta$, $y = r\sin\theta$ :
\[
J_\Phi = \begin{pmatrix}\cos\theta & -r\sin\theta\\\sin\theta & r\cos\theta\end{pmatrix} \qquad \det J_\Phi = r\cos^2\theta + r\sin^2\theta = r
\]
C'est le jacobien des coordonnées polaires, qui apparaît dans $\d x\,\d y = r\,\d r\,\d\theta$.

\textbf{4.} $G = (G_1,G_2)$ avec $G_1 = u^2-v$, $G_2 = uv$ :
\[
J_G(u,v) = \begin{pmatrix}2u & -1\\v & u\end{pmatrix}
\]
En $(u,v) = F(1,2) = (1+2, 1\times2+4) = (3, 6)$ :
\[
J_G(3,6) = \begin{pmatrix}6 & -1\\6 & 3\end{pmatrix}
\]
Par la règle de la chaîne : $J_{G\circ F}(1,2) = J_G(F(1,2))\cdot J_F(1,2)$ :
\[
J_{G\circ F}(1,2) = \begin{pmatrix}6 & -1\\6 & 3\end{pmatrix}\begin{pmatrix}2 & 1\\2 & 5\end{pmatrix} = \begin{pmatrix}12-2 & 6-5\\12+6 & 6+15\end{pmatrix} = \begin{pmatrix}10 & 1\\18 & 21\end{pmatrix}
\]
\end{boitecorrige}

% ----------------------
\begin{boiteexercice}
\textbf{Exercice 9 - Extrema libres et conditions du second ordre}\niveau{\moyen}

Trouver et classifier tous les points critiques de :
\[
f(x,y) = x^4 + y^4 - 4xy
\]
\end{boiteexercice}

\begin{boitecorrige}
\textbf{Points critiques.} On résout $\nabla f = \vec{0}$ :
\[
\dpart{f}{x} = 4x^3 - 4y = 0 \Rightarrow y = x^3
\]
\[
\dpart{f}{y} = 4y^3 - 4x = 0 \Rightarrow x = y^3
\]
En substituant : $y = x^3 = (y^3)^3 = y^9$, donc $y^9 - y = 0$, soit $y(y^8-1) = 0$.

Solutions : $y = 0$ ou $y^8 = 1$ (soit $y = \pm1$ dans $\mathbb{R}$).
\begin{itemize}[nosep]
  \item $y = 0 \Rightarrow x = 0$ : point $(0,0)$.
  \item $y = 1 \Rightarrow x = y^3 = 1$ : point $(1,1)$.
  \item $y = -1 \Rightarrow x = -1$ : point $(-1,-1)$.
\end{itemize}

\textbf{Matrice hessienne.}
\[
H_f(x,y) = \begin{pmatrix}12x^2 & -4\\-4 & 12y^2\end{pmatrix}
\]

En $(0,0)$ : $H_f = \begin{pmatrix}0&-4\\-4&0\end{pmatrix}$, $\det H = -16 < 0$ : \textbf{point selle}.

En $(1,1)$ : $H_f = \begin{pmatrix}12&-4\\-4&12\end{pmatrix}$, $\det H = 144-16 = 128 > 0$ et $H_{11} = 12 > 0$ : \textbf{minimum local}. $f(1,1) = 1+1-4 = -2$.

En $(-1,-1)$ : même matrice hessienne, même conclusion : \textbf{minimum local}.
\[
f(-1,-1) = 1+1-4(-1)(-1) = -2.
\]

Les deux minima locaux ont la même valeur $f = -2$, ils sont globaux par argument de croissance à l'infini.
\end{boitecorrige}

% ----------------------
\begin{boiteexercice}
\textbf{Exercice 10 - Intégrale double sur un domaine non rectangulaire}\niveau{\moyen}

Calculer l'intégrale $I = \displaystyle\iint_D (x+y)\,\mathrm{d}x\,\mathrm{d}y$ sur le domaine $D$ délimité par $y = x^2$ et $y = x$ (avec $0\leq x\leq1$).
\begin{enumerate}
  \item Tracer le domaine $D$ et écrire les bornes d'intégration.
  \item Calculer l'intégrale par Fubini.
  \item Vérifier en inversant l'ordre d'intégration.
\end{enumerate}
\end{boiteexercice}

\begin{boitecorrige}
\textbf{1.} Les courbes $y=x^2$ et $y=x$ se croisent en $x=0$ et $x=1$. Pour $0\leq x\leq1$, $x^2\leq x$, donc $D = \{(x,y) : 0\leq x\leq1,\; x^2\leq y\leq x\}$.

\textbf{2.} Par Fubini (intégrer $y$ en premier) :
\[
I = \int_0^1\!\int_{x^2}^x(x+y)\,\mathrm{d}y\,\mathrm{d}x = \int_0^1\!\left[xy + \frac{y^2}{2}\right]_{x^2}^x\mathrm{d}x
\]
\[
= \int_0^1\!\left[\left(x^2+\frac{x^2}{2}\right) - \left(x^3+\frac{x^4}{2}\right)\right]\mathrm{d}x = \int_0^1\!\left(\frac{3x^2}{2} - x^3 - \frac{x^4}{2}\right)\mathrm{d}x
\]
\[
= \left[\frac{x^3}{2} - \frac{x^4}{4} - \frac{x^5}{10}\right]_0^1 = \frac{1}{2} - \frac{1}{4} - \frac{1}{10} = \frac{10-5-2}{20} = \frac{3}{20}
\]

\textbf{3.} Inversons l'ordre : pour $0\leq y\leq1$, $x$ varie entre $y$ et $\sqrt{y}$ (car $x\geq y\Leftrightarrow x\geq y$ et $x\leq\sqrt{y}$) :
\[
I = \int_0^1\!\int_y^{\sqrt{y}}(x+y)\,\mathrm{d}x\,\mathrm{d}y = \int_0^1\!\left[\frac{x^2}{2}+yx\right]_y^{\sqrt{y}}\mathrm{d}y
\]
\[
= \int_0^1\!\left[\left(\frac{y}{2}+y\sqrt{y}\right)-\left(\frac{y^2}{2}+y^2\right)\right]\mathrm{d}y = \int_0^1\!\left(\frac{y}{2}+y^{3/2}-\frac{3y^2}{2}\right)\mathrm{d}y
\]
\[
= \left[\frac{y^2}{4}+\frac{2y^{5/2}}{5}-\frac{y^3}{2}\right]_0^1 = \frac{1}{4}+\frac{2}{5}-\frac{1}{2} = \frac{5+8-10}{20} = \frac{3}{20} \quad
\]
\end{boitecorrige}

% ----------------------
\begin{boiteexercice}
\textbf{Exercice 11 - Théorème des fonctions implicites}\niveau{\difficile}

Soit $F(x,y) = x^2 + y^2 - 1$ (cercle unité). On cherche à exprimer $y$ comme fonction de $x$ localement.
\begin{enumerate}
  \item En quel point $(x_0, y_0) = (0,1)$ le théorème des fonctions implicites s'applique-t-il ? Quelle condition vérifie $\partial F/\partial y$ ?
  \item Donner la formule de la dérivée $\d y/\d x$ implicitement définie par $F(x,y)=0$.
  \item Calculer $\d y/\d x$ en $(0,1)$ et en $(1/\sqrt{2}, 1/\sqrt{2})$.
  \item Calculer $\d^2y/\d x^2$ en différentiant à nouveau la relation implicite.
\end{enumerate}
\end{boiteexercice}

\begin{boitecorrige}
\textbf{1.} En $(0,1)$ : $F(0,1) = 0+1-1 = 0$. On vérifie $\dpart{F}{y}(0,1) = 2y\big|_{(0,1)} = 2\neq0$. Le théorème des fonctions implicites s'applique : il existe $\varphi$ définie au voisinage de 0 telle que $F(x,\varphi(x)) = 0$ et $\varphi(0) = 1$.

\textbf{2.} En différentiant $F(x,y(x)) = 0$ par rapport à $x$ :
\[
\dpart{F}{x} + \dpart{F}{y}\cdot\frac{\d y}{\d x} = 0 \quad\Rightarrow\quad \frac{\d y}{\d x} = -\frac{\partial F/\partial x}{\partial F/\partial y} = -\frac{2x}{2y} = -\frac{x}{y}
\]

\textbf{3.}
\begin{itemize}[nosep]
  \item En $(0,1)$ : $\d y/\d x = -0/1 = 0$ (tangente horizontale au sommet du cercle).
  \item En $(1/\sqrt{2}, 1/\sqrt{2})$ : $\d y/\d x = -(1/\sqrt{2})/(1/\sqrt{2}) = -1$ (tangente à $-45°$).
\end{itemize}

\textbf{4.} On dérive $\d y/\d x = -x/y$ par rapport à $x$ :
\[
\frac{\d^2y}{\d x^2} = -\frac{y - x(\d y/\d x)}{y^2} = -\frac{y - x(-x/y)}{y^2} = -\frac{y^2+x^2}{y^3} = -\frac{1}{y^3}
\]
(car $x^2+y^2=1$). En $(0,1)$ : $\d^2y/\d x^2 = -1$ : le cercle est concave vers le bas, courbure $= 1$.
\end{boitecorrige}

\subsection*{Exercice 12 - Ligne de niveau et domaine de définition \niveau{\moyen}}

\begin{boiteexercice}
Soit la fonction $f(x,y) = \ln\!\left(x^2 - 6x + y^2 + 4y - 2\right)$.
\begin{enumerate}
  \item Déterminer le domaine de définition $\mathcal{D}_f$ de $f$ et le représenter géométriquement.
  \item Déterminer la ligne de niveau $z = 0$ de $f$, c'est-à-dire l'ensemble $\{(x,y)\in\mathcal{D}_f : f(x,y) = 0\}$.
\end{enumerate}
\end{boiteexercice}

\begin{boitecorrige}
\textbf{1. Domaine de définition.} La fonction $\ln$ est définie pour un argument strictement positif :
\[
x^2 - 6x + y^2 + 4y - 2 > 0.
\]
On complète les carrés :
\[
x^2 - 6x + y^2 + 4y - 2 = (x-3)^2 - 9 + (y+2)^2 - 4 - 2 = (x-3)^2 + (y+2)^2 - 15.
\]
La condition devient $(x-3)^2 + (y+2)^2 > 15$. Le domaine $\mathcal{D}_f$ est donc le plan $\R^2$ privé du disque fermé de centre $\Omega(3,-2)$ et de rayon $R = \sqrt{15}$.

\textbf{2. Ligne de niveau $z=0$.} On résout $f(x,y) = 0$, soit :
\[
\ln\!\left((x-3)^2 + (y+2)^2 - 15\right) = 0 \iff (x-3)^2 + (y+2)^2 - 15 = 1.
\]
D'où :
\[
\boxed{(x-3)^2 + (y+2)^2 = 16}
\]
qui est l'équation du cercle de centre $\Omega(3,-2)$ et de rayon $4$. Puisque $4 > \sqrt{15}$, ce cercle est bien contenu dans $\mathcal{D}_f$ : la ligne de niveau est donc le cercle complet.

\reponse{La ligne de niveau $z=0$ est le cercle $(x-3)^2 + (y+2)^2 = 16$, de centre $(3,-2)$ et de rayon $4$.}
\end{boitecorrige}

\subsection*{Exercice 13 - Calcul de limites au point $(0,0)$ \niveau{\moyen}}

\begin{boiteexercice}
Calculer, si elle existe, la limite en $(0,0)$ des fonctions suivantes :
\begin{enumerate}
  \item $f_1(x,y) = \dfrac{(1+x+y)\sin y}{y}$.
  \item $f_2(x,y) = (x^2+y^2)\sin\!\left(\dfrac{\pi}{x^2+y^2}\right)$.
  \item $f_3(x,y) = \dfrac{xy}{x^2+y^2}$.
\end{enumerate}
\end{boiteexercice}

\begin{boitecorrige}
\textbf{1.} On réécrit $f_1(x,y) = (1+x+y)\cdot\dfrac{\sin y}{y}$. Comme $\dfrac{\sin y}{y} \xrightarrow[y\to 0]{} 1$ et $1+x+y \xrightarrow[(x,y)\to(0,0)]{} 1$, on a par produit :
\[
\lim_{(x,y)\to(0,0)} f_1(x,y) = 1 \times 1 = \boxed{1}.
\]

\textbf{2.} On utilise le théorème des gendarmes. Comme $\left|\sin\!\left(\dfrac{\pi}{x^2+y^2}\right)\right| \leq 1$, on a :
\[
0 \leq \left|(x^2+y^2)\sin\!\left(\dfrac{\pi}{x^2+y^2}\right)\right| \leq x^2+y^2 \xrightarrow[(x,y)\to(0,0)]{} 0.
\]
Par gendarmes : $\lim_{(x,y)\to(0,0)} f_2(x,y) = \boxed{0}$.

\textbf{3.} Cette limite n'existe pas. En approchant selon la droite $y = mx$ :
\[
f_3(x, mx) = \dfrac{x(mx)}{x^2 + (mx)^2} = \dfrac{m}{1+m^2},
\]
qui dépend de $m$. Par exemple, selon $y=0$ ($m=0$) on obtient $0$, et selon $y=x$ ($m=1$) on obtient $\dfrac{1}{2}$. La limite dépend donc du chemin d'approche, elle n'existe pas.

\reponse{$f_1 \to 1$, $f_2 \to 0$, $f_3$ n'a pas de limite en $(0,0)$.}
\end{boitecorrige}

\subsection*{Exercice 14 - Optimisation : durée d'une infection \niveau{\difficile}}

\begin{boiteexercice}
Une étude en laboratoire montre que la durée d'une infection bactérienne (en jours) peut être modélisée par la fonction
\[
D(x, y) = x^2 + 2y^2 - 18x - 24y + 2xy + 120,
\]
où $x$ et $y$ sont les doses (en mg) de deux composés chimiques.
\begin{enumerate}
  \item Calculer les dérivées partielles $\dpart{D}{x}$ et $\dpart{D}{y}$.
  \item Déterminer le point critique $(x^\star, y^\star)$ de $D$.
  \item Étudier la nature du point critique à l'aide de la matrice hessienne.
  \item En déduire les doses optimales et la durée minimale de l'infection.
\end{enumerate}
\end{boiteexercice}

\begin{boitecorrige}
\textbf{1. Dérivées partielles.}
\[
\dpart{D}{x} = 2x - 18 + 2y, \qquad \dpart{D}{y} = 4y - 24 + 2x.
\]

\textbf{2. Point critique.} On résout $\dpart{D}{x} = \dpart{D}{y} = 0$ :
\[
\begin{cases}
2x + 2y = 18 \;\Rightarrow\; x + y = 9\\
2x + 4y = 24 \;\Rightarrow\; x + 2y = 12
\end{cases}
\]
Par soustraction : $y = 3$, puis $x = 6$. Le point critique est $(x^\star, y^\star) = (6, 3)$.

\textbf{3. Matrice hessienne.}
\[
H(x,y) = \begin{pmatrix} \dpart{^2 D}{\partial x^2} & \dpart{^2 D}{\partial x\partial y}\\[4pt] \dpart{^2 D}{\partial y\partial x} & \dpart{^2 D}{\partial y^2}\end{pmatrix} = \begin{pmatrix} 2 & 2 \\ 2 & 4 \end{pmatrix}.
\]
En $(6,3)$ : $\det H = 2\times 4 - 2\times 2 = 4 > 0$ et $\dpart{^2 D}{\partial x^2} = 2 > 0$. Donc $H$ est définie positive : $(6,3)$ est un \textbf{minimum local} (et global, car $D$ est quadratique strictement convexe).

\textbf{4. Doses optimales et durée minimale.}
\[
D(6,3) = 36 + 18 - 108 - 72 + 36 + 120 = 30 \text{ jours.}
\]

\reponse{Doses optimales : $x^\star = 6$ mg, $y^\star = 3$ mg ; durée minimale $D_{\min} = 30$ jours.}
\end{boitecorrige}

\ifdefined\inMainDoc\else\end{document}\fi
